
\documentclass[a4paper]{article}
\usepackage[T1]{fontenc}
\usepackage{amssymb,amsmath}
\usepackage{multicol,multirow}
\usepackage{array}
\usepackage{calc}
\usepackage{ifthen}
\usepackage{titlesec}
\usepackage{enumitem}
\usepackage[landscape]{geometry}
\usepackage[colorlinks=true,citecolor=blue,linkcolor=blue]{hyperref}
\usepackage{xcolor}

\setcounter{secnumdepth}{4}

\titleformat{\paragraph}
{\normalfont\normalsize\bfseries}{\theparagraph}{1em}{}
\titlespacing*{\paragraph}
{0pt}{3.25ex plus 1ex minus .2ex}{1.5ex plus .2ex}

\ifthenelse{\lengthtest { \paperwidth = 11in}}
    { \geometry{top=.5in,left=.5in,right=.5in,bottom=.5in} }
	{\ifthenelse{ \lengthtest{ \paperwidth = 297mm}}
		{\geometry{top=1cm,left=1cm,right=1cm,bottom=1cm} }
		{\geometry{top=1cm,left=1cm,right=1cm,bottom=1cm} }
	}

\pagestyle{empty}

% Section: centered in a light gray box with black outline
\newcommand{\sectionbox}[1]{%
  \fcolorbox{black}{black!8}{\parbox{\dimexpr\columnwidth-2\fboxsep-2\fboxrule\relax}{\centering #1}}}
\titleformat{\section}[block]
  {\normalfont\large\bfseries}{}{0mm}{\sectionbox}
\titlespacing*{\section}{0mm}{0.3ex plus .3ex minus .2ex}{0.3ex plus .2ex}

\makeatletter
\renewcommand{\subsection}{\@startsection{subsection}{2}{0mm}%
                                {-0.8ex plus -.5ex minus -.2ex}%
                                {0.3ex plus .2ex}%
                                {\normalfont\normalsize\bfseries}}
\renewcommand{\subsubsection}{\@startsection{subsubsection}{3}{0mm}%
                                {-0.8ex plus -.5ex minus -.2ex}%
                                {0.5ex plus .2ex}%
                                {\normalfont\small\bfseries}}
\makeatother
\setcounter{secnumdepth}{0}
\setlength{\parindent}{0pt}
\setlength{\parskip}{0pt plus 0.3ex}

% -----------------------------------------------------------------------
% Keybinding notation
% -----------------------------------------------------------------------
\definecolor{keybg}{RGB}{234,234,244}
\definecolor{warn}{RGB}{155,25,25}

% \key{SPC g g} -- a chord or sequence, shaded so it reads as a key.
% \fboxsep is trimmed so a key box fits inside a narrow table column.
\newcommand{\key}[1]{{\setlength{\fboxsep}{1.2pt}%
  \colorbox{keybg}{\strut\texttt{\footnotesize #1}}}}
% \danger{...} -- destructive commands
\newcommand{\danger}[1]{\textcolor{warn}{#1}}
% \fn{magit-status} -- the command a key runs
\newcommand{\fn}[1]{\texttt{\footnotesize #1}}

% Two-column key table: keys on the left, what they do on the right.
% Widths are computed from \columnwidth rather than using tabularx, because
% tabularx scans for a literal \end{tabularx} and cannot be wrapped like this.
\newenvironment{keys}[1][2.9cm]
  {\par\nopagebreak\begin{tabular}%
     {@{}>{\raggedright\arraybackslash}p{#1}@{\hspace{4pt}}%
       >{\raggedright\arraybackslash}p{\dimexpr\columnwidth-#1-4pt\relax}@{}}}
  {\end{tabular}\par\vspace{0.4ex}}

% One row of a key table.
\newcommand{\krow}[2]{\key{#1} & #2 \\}

% A row whose key cell is marked up by hand, for alternatives that need to be
% separate boxes so the column can break the line between them.
\newcommand{\krowx}[2]{#1 & #2 \\}
\newcommand{\korr}[2]{\key{#1}\,/\,\allowbreak\key{#2}}

% Escapes for keys that collide with LaTeX syntax
\newcommand{\ktilde}{\textasciitilde}
\newcommand{\kcaret}{\textasciicircum}
\newcommand{\kbslash}{\textbackslash}
